% Section 4b — L1 Agentic Parkinson, measured as a budget elasticity

\section{The Wall Axis: L1 Measured, Not Assumed}
\label{sec:l1}

Of the three laws, L1 is the one we could most easily have left as an
assertion. Parkinson's observation that ``work expands so as to fill the time
available'' \citep{parkinson1955} is the folk theory of what a budgeted agent
does, and an earlier draft of this paper stated its contrapositive --- that
native agents simply do not exhibit it --- on the strength of the $\CAR$
numbers in \S\ref{sec:l2}. That statement was wrong, and the corpus already
contained the data to show it.

\paragraph{Why not T1.3.} T1.3 is the sub-capability nominally assigned to
L1, and it cannot answer the question. Its trajectories carry
\texttt{budget\_kind="none"}, so the stated deadline lives only inside the
prompt string and $\parkinson$ is not computable from them. More seriously,
its ten prompts pair each deadline with a different question, and the pairing
tracks intrinsic task size almost perfectly: the five-second item asks for the
capital of Australia, the 180-second item asks for a short story. Deadline and
task size are confounded, so a longer answer under a longer deadline is
uninformative. No reanalysis repairs this; the design has to be re-run
crossed, which we describe in \S\ref{sec:limitations}.

\paragraph{The design L1 needs is already in T2.3.} The wall-budget
capability holds the task fixed and varies the budget over
$\budget \in \{60, 300, 900, 1800, 3600\}$\,s at $n = 30$ per cell per
setting --- a $60\times$ range on a single factor. We therefore measure L1 as
a \emph{budget elasticity},
\[
  e \;:=\; \frac{\mathrm{d}\,\log_{10} \twallstar}{\mathrm{d}\,\log_{10} \budget},
\]
fit per agent by least squares over all $(\budget, \twallstar)$ pairs with a
$95\%$ bootstrap CI ($4{,}000$ resamples). Parkinson's law in its strong form
is $e = 1$: the work expands to fill whatever it is given. Complete
indifference to the budget is $e = 0$.

\begin{table}[h]
\centering\small
\caption{L1 budget elasticity across the panel. $e = 1$ is Parkinson's law in
full; $e = 0$ is no response to the budget at all. Pooled over both settings,
$n = 60$ per agent.}
\label{tab:l1-elasticity}
\begin{tabular}{lcc|lcc}
\toprule
Agent & $e$ & 95\% CI & Agent & $e$ & 95\% CI \\
\midrule
Claude Haiku 4.5 & $+0.228$ & $[+0.170, +0.284]$ & MiniMax-M2.7   & $+0.089$ & $[+0.028, +0.147]$ \\
GLM-5.2-FP8      & $+0.219$ & $[+0.166, +0.272]$ & gpt-4o-mini    & $+0.082$ & $[-0.141, +0.219]$ \\
gpt-5.1          & $+0.149$ & $[+0.102, +0.202]$ & Sonnet $+$ thinking & $+0.074$ & $[-0.038, +0.185]$ \\
o3               & $+0.127$ & $[+0.062, +0.196]$ & o4-mini        & $+0.072$ & $[+0.028, +0.131]$ \\
gpt-4o           & $+0.103$ & $[+0.068, +0.139]$ & Qwen3.6-27B    & $+0.058$ & $[+0.036, +0.079]$ \\
Kimi-K2.6        & $+0.092$ & $[+0.027, +0.156]$ & Claude Sonnet 4.6 & $+0.052$ & $[+0.021, +0.090]$ \\
                 &          &                    & Qwen2.5-7B     & $-0.020$ & $[-0.099, +0.061]$ \\
\midrule
\multicolumn{6}{l}{\textbf{Panel median $e = 0.089$}; 10 of 13 agents have a CI strictly above zero.} \\
\bottomrule
\end{tabular}
\end{table}

\paragraph{Result: present, and an order of magnitude too weak.}
The panel median is $e = 0.089$, and ten of thirteen agents across seven
vendors have a bootstrap CI strictly above zero (Table~\ref{tab:l1-elasticity}).
A $60\times$ increase in the stated budget therefore buys $1.4\times$ more
wall-clock, where Parkinson's law requires $60\times$. A single exponent
describes each agent's ladder to between $3\%$ and $22\%$ (median ${\sim}12\%$),
so the power-law summary is a reasonable description of the budget response and
not a line drawn through noise.

This is the narrative/action split of \S\ref{sec:l2}--\ref{sec:l3} visible on a
single axis. The agent's response to the budget is \emph{directionally correct}
--- it reads the number and moves the right way --- and it is roughly a tenth of
the strength the instruction calls for. An agent that could not read the budget
at all would give $e = 0$; an agent that could act on what it read would give
$e = 1$. The panel sits an order of magnitude from the second and, for most
agents, reliably away from the first.

\paragraph{What this is not.} Two limits bound the claim, and both point at the
same missing experiment. First, T2.3 is a \emph{single task}: $e$ is measured on
one item, so we cannot say whether the elasticity is a property of the agent or
of that item. Second, the T2.3 prompt explicitly instructs the agent to spend
the budget (``\textit{work on this task for $\budget$ seconds; do not stop
early}''), so $e$ measures compliance with an instruction, not the spontaneous
expansion Parkinson described. That makes $e = 0.089$ an \emph{upper bound} on
spontaneous L1 rather than a measurement of it --- which strengthens the
direction of the finding while weakening its interpretation as Parkinson's law
proper. The crossed design that would fix both, questions $\times$ deadlines
fully crossed with the budget recorded as a field, is specified in
\S\ref{sec:limitations}.

\paragraph{Relation to $\CAR$.} Since $\CAR = \twallstar/\budget$, we have
$\mathrm{d}\log \CAR / \mathrm{d}\log \budget = e - 1$ identically. The steep
decay of $\CAR$ with budget reported in \S\ref{sec:l2} is therefore the same
measurement in different coordinates, not independent corroboration of it, and
we do not present it as such. What the two views jointly establish is that one
exponent near zero governs both: the agent's wall-clock expenditure is nearly
invariant to what it is told it has.
